% This is a modified version of Springer's LNCS template suitable for anonymized MICCAI 2025 main conference submissions. 
% Original file: samplepaper.tex, a sample chapter demonstrating the LLNCS macro package for Springer Computer Science proceedings; Version 2.21 of 2022/01/12

\documentclass[runningheads]{llncs}
%
\usepackage[T1]{fontenc}
% T1 fonts will be used to generate the final print and online PDFs,
% so please use T1 fonts in your manuscript whenever possible.
% Other font encodings may result in incorrect characters.
%
\usepackage{graphicx,verbatim}
\usepackage{amsmath}
\usepackage{float}
\usepackage{booktabs}

% Used for displaying a sample figure. If possible, figure files should
% be included in EPS format.
%
% If you use the hyperref package, please uncomment the following two lines
% to display URLs in blue roman font according to Springer's eBook style:
%\usepackage{color}
%\renewcommand\UrlFont{\color{blue}\rmfamily}
%\urlstyle{rm}
%
\begin{document}
%
\title{On the Cone Effect and Modality Misalignment in Medical Vision–Language Embeddings}
%\titlerunning{Abbreviated paper title}
% If the paper title is too long for the running head, you can set
% an abbreviated paper title here
%
\begin{comment}  %% Removed for anonymized MICCAI submission
\author{First Author\inst{1}\orcidID{0000-1111-2222-3333} \and
Second Author\inst{2,3}\orcidID{1111-2222-3333-4444} \and
Third Author\inst{3}\orcidID{2222--3333-4444-5555}}
%
\authorrunning{F. Author et al.}
% First names are abbreviated in the running head.
% If there are more than two authors, 'et al.' is used.
%
\institute{Princeton University, Princeton NJ 08544, USA \and
Springer Heidelberg, Tiergartenstr. 17, 69121 Heidelberg, Germany
\email{lncs@springer.com}\\
\url{http://www.springer.com/gp/computer-science/lncs} \and
ABC Institute, Rupert-Karls-University Heidelberg, Heidelberg, Germany\\
\email{\{abc,lncs\}@uni-heidelberg.de}}

\end{comment}

\author{Anonymized Authors}  %% Added for anonymized MICCAI submission
\authorrunning{Anonymized Author et al.}
\institute{Anonymized Affiliations \\
    \email{email@anonymized.com}}
  
\maketitle              % typeset the header of the contribution
%
\begin{abstract}
The abstract should briefly summarize the contents of the paper in 150--250 words.  If you are to include a link to your Repository, please make sure it is anonymized for the double-blind review phase.

\keywords{First keyword  \and Second keyword \and Another keyword.}
% Authors must provide keywords and are not allowed to remove this Keyword section.

\end{abstract}
%
%
%
\section{Introduction}
Self-supervised learning via contrastive language-image pretraining (CLIP) has established itself as a cornerstone for developing frontier solutions in both natural and medical imaging domains \cite{radford2021learning,zhang2022contrastive}. By leveraging multi-modal supervision, this framework allows to capture complementary information from both modalities in a joint representation space. However, the contraction mapping enforced by the composition of non-linearities in deep networks inevitably leads to the so called ``cone effect'', whereby representations are mapped to a very dense and compact region in representation space. As CLIP objectives rely on separate encoders for image and text, this phenomenon causes embeddings to cluster in narrow, modality-specific regions of the representation manifold. These clusters can remain arbitrarily distant, a phenomenon termed \emph{modality gap} \cite{liang2022mind}. 

It is still not clear if the existence of a domain gap is pathological, although it is known from its inception that it drastically affects downstream performance \cite{liang2022mind}. Empirical and theoretical literature also suggest that CLIP alignment is also suboptimal in terms of intra and inter modality alignment \cite{liang2022mind,shi2023towards,zhang2023diagnosing,schrodi2024two,lee2025diffusion,mistretta2025cross}.

The simplest mechanism of action to correct potential misalignment is the length of the domain-gap manipulate in the embedding space. However, from an information theoretic perspective, simply narrowing the gap as much as possible may lead to perfect inter-modality alignment, but comes at a potentially high information cost exclusive to each modality \cite{jiang2023understanding}. It is possible to minimize this gap in a more constrained fashion \cite{shi2023towards,hu2024reclip,eslami2024mitigate}, but current methods need to be integrated during the pre-training stage of the CLIP model itself. This is also the case for other approaches that address inter or intra-class alignment, relying on explicit terms in the pre-text objective \cite{jiang2023understanding,mistretta2025cross} or training that depends on model inversion methods \cite{huang2025mind,lee2025diffusion}.  

In practice, CLIP-based models are rarely trained from scratch; they are fine-tuned, adapted, or probed. Consequently, there is a critical need for lightweight mechanisms that can manipulate the modality gap post-hoc, adapting pre-trained representations to downstream requirements without expensive retraining. We thus propose a method that takes advantage of the cone effect -- since each modality will be supported in a very dense and compact set in the representation space, a single learnable parameter $\lambda$ can be used to displace the intra-modality embeddings around their centroids, increasing/decreasing the domain gap depending on the downstream task. 

Our contributions are summarized as follows:
\begin{enumerate}
\item We introduce a framework that keeps the CLIP modality embedders frozen. The optimal downstream domain gap is modeled via a single learnable parameter in the downstream task.
\item We perform an extensive analysis on a collection of VLMs pre-trained on both medical and natural images, across a variety of datasets, both comprised of natural or medical images.
\item We study the impact of the proposed data-driven domain gap adaptation in both early and late fusion settings.
\end{enumerate}

%\section{Notes and todos}
%\begin{enumerate}
%    \item Add more refs in terms of the success of these types of VLMs in solving tasks in the intro/SOTA
%    \item If we have time the whitening transform could be cool in the fusion stage \cite{zhang2024dual}: idea extract orthogonal directions and vary the gap only on those directions
%    \item See if pretty pictures of \cite{yaras2024explaining} make sense
%\end{enumerate}


% =========================
% 2.4.3 / 2.4.4 / 2.4.5 (REPLACEMENT TEXT)
% Drop-in replacement for your current subsections under:
% \subsubsection{Reducing the Modality Embedding Gap}
% =========================

\section{Methods}
\label{sec:methods}

\subsection{Problem Statement: Cone Effect and Modality Gap in Medical VLMs}
\label{sec:problem_theory}

\subsubsection{Modality Gap in Vision--Language Embedding Spaces}
\label{sec:modality_gap}

Vision--Language Models (VLMs) such as CLIP and SigLIP learn joint embedding spaces by mapping images and texts into a shared hyperspherical representation. Let $f_\theta:\mathcal{I}\rightarrow\mathbb{R}^d$ and $g_\phi:\mathcal{T}\rightarrow\mathbb{R}^d$ denote the image and text encoders. We consider normalized embeddings

\begin{equation}
v=\frac{f_\theta(I)}{\|f_\theta(I)\|_2},
\qquad
t=\frac{g_\phi(T)}{\|g_\phi(T)\|_2},
\qquad
v,t\in\mathbb{S}^{d-1}.
\end{equation}

Despite contrastive supervision, embeddings from each modality remain separated by a persistent displacement, termed the \emph{modality gap}~\cite{liang2022mind}. We quantify this gap through the centroid difference

\begin{equation}
\mu_v=\mathbb{E}[v],
\qquad
\mu_t=\mathbb{E}[t],
\qquad
\Delta=\mu_v-\mu_t,
\label{eq:gap_vector}
\end{equation}
where $\|\Delta\|_2$ measures global cross-modal separation.

% ------------------------------------------------------------

\subsubsection{Cone Effect and Modality Gap}
\label{sec:cone_effect}

A central geometric driver of the modality gap is the \emph{cone effect}, which describes the tendency of deep embeddings to concentrate within narrow angular regions of the hypersphere~\cite{liang2022mind}. In ReLU networks, activations restrict representations to positive orthants, $\phi(x)=\max(x,0)$, and Liang et al.~\cite{liang2022mind} show that cosine similarity increases monotonically through depth,
\begin{equation}
\cos(\phi(Wx_1),\phi(Wx_2)) \geq \cos(x_1,x_2),
\label{eq:cone_similarity}
\end{equation}
inducing anisotropic embedding cones around dominant mean directions.
To quantify the degree of angular concentration, we use the Mean Resultant Length, $R = \|\mathbb{E}[v]\|_2$. For embeddings on the unit hypersphere, $R \in [0, 1]$, where values closer to 1 indicate highly concentrated distributions (narrow cones) and values near 0 indicate isotropy.

Contrastive objectives further stabilize this geometry through the interaction of attractive and repulsive forces. 
In CLIP, the InfoNCE loss can be decomposed as:

\begin{equation}
-\log\frac{\exp(v_i^\top t_i/\tau)}
{\sum_{j}\exp(v_i^\top t_j/\tau)}
=
\underbrace{-\frac{1}{\tau}v_i^\top t_i}_{\text{attractive alignment}}
+
\underbrace{\log\sum_{j}\exp(v_i^\top t_j/\tau)}_{\text{repulsive separation}},
\label{eq:clip_attract_repel}
\end{equation}


where the first term is an \emph{attractive alignment} term pulling matched pairs together, while the second term is a \emph{batch-wise repulsive} term pushing embeddings away from all negatives.

SigLIP replaces softmax normalization with an all-pairs logistic formulation,
\begin{equation}
\mathcal{L}_{\text{SigLIP}}
=
-\sum_{i,j}\log\sigma\!\big(L_{ij}(\alpha v_i^\top t_j+\beta)\big),
\end{equation}
where negatives contribute independently and gradients saturate once pairs become sufficiently separated. 
Consequently, the effective strength of repulsion depends on the negative-to-positive ratio and the number of hard negatives within the batch, which can preserve modality-specific cones more strongly under limited semantic diversity.



\subsubsection{Cone Concentration in Medical Domains}
\label{sec:medical_specifics}

These effects are amplified in medical domains, where the medical data distribution $\mathcal{M}$ constitutes a low-variance subset of the natural distribution $\mathcal{D}$. Represented as ($\mathcal{M}\subseteq\mathcal{D}$). 
To formalize this, we adopt the variance decomposition of intermediate layer outputs introduced by Liang et al.~\cite{liang2022mind}. 
Let $U$ denote the input to a given encoder layer and let $h_\Theta(U)$ denote its output under the layer parameters $\Theta$. 

Then the total output variance satisfies the law of total variance:
\begin{equation}
\mathrm{Var}\!\left[h_\Theta(U)\right]
=
\underbrace{\mathbb{E}_\Theta\!\left[\mathrm{Var}\!\left(h_\Theta(U)\mid\Theta\right)\right]}_{\text{data-induced variance}}
+
\underbrace{\mathrm{Var}_\Theta\!\left(\mathbb{E}\!\left[h_\Theta(U)\mid\Theta\right]\right)}_{\text{weight-induced variance}}.
\label{eq:total_variance_liang}
\end{equation}

In medical datasets, reduced semantic and visual diversity decreases the varianze in the data-induced component,
\begin{equation}
\mathbb{E}_\Theta\!\left[\mathrm{Var}\!\left(h_\Theta(U_{\text{med}})\mid\Theta\right)\right]
<
\mathbb{E}_\Theta\!\left[\mathrm{Var}\!\left(h_\Theta(U_{\text{nat}})\mid\Theta\right)\right],
\label{eq:data_variance_medical}
\end{equation}
making medical representations even more susceptible to the cone effect. 

%On the hypersphere, this corresponds to lower embedding dispersion,
%\mathbb{E}[\|v-\mu_v\|^2]_{\text{med}} < \mathbb{E}[\|v-\mu_v\|^2]_{\text{nat}}$,
%implying stronger cone effects and potentially larger centroid displacement between modalities.

However, medical-tuned backbones introduce domain-specific inductive biases through specialized pretraining. Such adaptation modifies the weight-induced term in Eq.~\ref{eq:total_variance_liang} by reshaping the representation geometry, which can partially compensate for reduced data variance by expanding clinically meaningful directions. 

%As a result, medical pretraining produces qualitatively different cone configurations and modality gap behavior across medical datasets.


\subsubsection{Hyperspherical Embedding Alignment Intervention}
\label{sec:alignment}

To mitigate the modality gap without modifying encoder weights, we apply an inference-time embedding alignment intervention. For $\lambda\in[0,1]$, embeddings are shifted along the gap vector and renormalized:
\begin{equation}
t'=\mathrm{Norm}\!\left(t+\frac{\lambda}{2}\Delta\right),
\qquad
v'=\mathrm{Norm}\!\left(v-\frac{\lambda}{2}\Delta\right).
\label{eq:alignment}
\end{equation}
By varying $\lambda$, we study controlled closure of $\|\Delta\|_2$ while preserving intra-modality cone geometry.


% ============================================================
\subsection{Experimental Setup}
\label{sec:experimental_setup}

\subsubsection{Vision--Language Backbones}
\label{sec:backbones}

To characterize the interaction between cone concentration, modality gap, and domain-specific inductive biases, we evaluate both general-domain and medically specialized dual-encoder VLMs. Specifically, we consider CLIP (ViT-B/32) and SigLIP (ViT-B/16) as natural-domain baselines, and BioMedCLIP and MedSigLIP as medical-domain variants trained on biomedical corpora and clinically curated image--text pairs. For each backbone, we extract frozen image and text embeddings independently using Eq.~(1), enabling direct geometric comparison across modalities and domains without confounding effects from downstream fine-tuning.

% ------------------------------------------------------------

\subsubsection{Datasets and Domain Shift Regimes}
\label{sec:datasets}

We conduct experiments on a four medical datasets covering different modalities and data types such as MIMIC-CXR, HAM10000, BRSET, and mBRSET. For reference, we also included some natural domain datasets, including Recipes5k, DAQUAR, COCO-QA, and Fakeddit. Each dataset $\mathcal{D}_k=\{(I_i,T_i,y_i)\}_{i=1}^{N_k}$ consists of paired image--text samples with labels $y_i$ corresponding to either multi-class or multi-label prediction tasks. 

All datasets are partitioned into disjoint training, validation, and test splits,
$\mathcal{D}_k=\mathcal{D}_k^{\text{train}}\cup\mathcal{D}_k^{\text{val}}\cup\mathcal{D}_k^{\text{test}}$,
and all alignment operators are computed exclusively on the training split to avoid test leakage.

\subsubsection{Embedding Alignment Protocol}
\label{sec:lambda_protocol}

To evaluate the causal role of modality gap reduction, we apply the hyperspherical alignment operator defined in Eq.~\ref{eq:alignment} with alignment strength $\lambda\in[0,1]$. For each dataset and backbone, the modality gap vector $\Delta$ is estimated from training embeddings as
\begin{equation}
\Delta_k = \mu_v^{(k)} - \mu_t^{(k)},
\qquad
\mu_v^{(k)}=\mathbb{E}_{(I,T)\sim \mathcal{D}_k^{\text{train}}}[v],
\quad
\mu_t^{(k)}=\mathbb{E}_{(I,T)\sim \mathcal{D}_k^{\text{train}}}[t].
\end{equation}

We then generate aligned embeddings $\{(v'_i,t'_i)\}$ for $\lambda\in\{0.0,0.1,\dots,1.0\}$, producing a controlled family of representation spaces in which $\|\Delta_k\|_2$ is progressively reduced while preserving intra-modality cone geometry.

% ------------------------------------------------------------

\subsubsection{Linear Probe Evaluation of Representation Quality}
\label{sec:linear_probe_eval}

To isolate embedding quality from nonlinear classifier capacity, we evaluate aligned representations using an Early Fusion Linear Probe. For each aligned pair $(v'_i,t'_i)$, we form the concatenated feature vector
\begin{equation}
x_i = [v'_i;t'_i]\in\mathbb{R}^{2d},
\end{equation}
and train a single affine classifier
\begin{equation}
\hat{y}_i = Wx_i+b,
\label{eq:probe_classifier}
\end{equation}
with no hidden layers or nonlinear activations.

Probe parameters $(W,b)$ are optimized on $\mathcal{D}_k^{\text{train}}$ by minimizing the empirical risk
\begin{equation}
\min_{W,b}\;
\frac{1}{|\mathcal{D}_k^{\text{train}}|}
\sum_{(x_i,y_i)\in\mathcal{D}_k^{\text{train}}}
\ell(\hat{y}_i,y_i),
\label{eq:erm_probe}
\end{equation}
where $\ell$ corresponds to class-weighted binary cross-entropy for multi-label datasets (e.g., MIMIC-CXR) and categorical cross-entropy for multi-class datasets (e.g., HAM10000). 

Optimization is performed using stochastic gradient descent with momentum $0.9$ and learning rate $3\times 10^{-4}$, with early stopping based on validation loss.

% ------------------------------------------------------------

\subsubsection{Evaluation Metrics and Statistical Protocol}
\label{sec:metrics_protocol}

Model performance is assessed on held-out test splits $\mathcal{D}_k^{\text{test}}$ using the Area Under the Receiver Operating Characteristic Curve (AUC), providing a threshold-independent measure of discriminative quality. For each backbone, dataset, and alignment strength $\lambda$, experiments are repeated across five independent random seeds. Final results are reported as mean $\pm$ standard deviation,
\begin{equation}
\mathrm{AUC}_{k,\lambda} =
\frac{1}{5}\sum_{s=1}^5 \mathrm{AUC}_{k,\lambda}^{(s)}
\;\;\pm\;\;
\mathrm{Std}\big(\mathrm{AUC}_{k,\lambda}^{(s)}\big),
\end{equation}
ensuring robustness to initialization and data-ordering effects.

% ------------------------------------------------------------

\subsubsection{Geometric Visualization of Cone Structure}
\label{sec:geometry_vis}

To complement quantitative evaluation, we visualize modality cones and centroid displacement using Principal Component Analysis (PCA). Normalized embeddings are projected into low-dimensional subspaces, allowing qualitative inspection of (i) cross-modal separation induced by $\Delta_k$, and (ii) intra-modality angular concentration associated with the cone effect.












\section{Results and Discussion}
\label{sec:results_discussion}

\begin{figure}[H]
    \centering
    \includegraphics[width=.8\linewidth]{Figures/Fig1_3d_4.pdf}
    \caption{Low-dimensional visualization of image and text embeddings across datasets and backbones. Projections are shown using the principal components. Image and text embeddings form modality-specific cones with strong angular concentration. A persistent centroid displacement is observed across all models, with tighter cones and reduced dispersion in medical datasets. Alignment radii $R$ quantify intra-modality concentration.}
    \label{fig:geometry}
\end{figure}


\begin{table}[h]
\centering
\caption{Mean Resultant Length ($R$), quantifying the degree of cone concentration for Text and Image embeddings. Higher values indicate narrower cones (stronger concentration). $R \in [0, 1]$.}
\resizebox{\textwidth}{!}{%
\begin{tabular}{lcccc c cccc}
\toprule
 & \multicolumn{4}{c}{\textbf{Text Embeddings ($R$)}} & & \multicolumn{4}{c}{\textbf{Image Embeddings ($R$)}} \\
\cmidrule{2-5} \cmidrule{7-10}
\textbf{Dataset} & \textbf{CLIP} & \textbf{SigLIP} & \textbf{BioMed} & \textbf{MedSig} & & \textbf{CLIP} & \textbf{SigLIP} & \textbf{BioMed} & \textbf{MedSig} \\
\midrule
\multicolumn{10}{l}{\textit{Medical Datasets}} \\
\midrule
MIMIC-CXR & 0.915 & 0.862 & 0.899 & 0.726 & & 0.970 & 0.974 & 0.910 & 0.735 \\
HAM10000 & 0.919 & 0.809 & 0.889 & 0.970 & & 0.920 & 0.933 & 0.941 & 0.858 \\
BRSET & 0.985 & 0.956 & 0.990 & 0.980 & & 0.986 & 0.980 & 0.964 & 0.731 \\
mBRSET & 0.990 & 0.973 & 0.975 & 0.957 & & 0.986 & 0.982 & 0.958 & 0.770 \\
\midrule
\multicolumn{10}{l}{\textit{Natural Datasets}} \\
\midrule
COCO-QA & 0.852 & 0.743 & 0.769 & 0.751 & & 0.691 & 0.702 & 0.739 & 0.638 \\
Fakeddit & 0.745 & 0.730 & 0.734 & 0.604 & & 0.654 & 0.657 & 0.702 & 0.597 \\
Recipes5k & 0.775 & 0.760 & 0.869 & 0.632 & & 0.810 & 0.791 & 0.807 & 0.738 \\
DAQUAR & 0.903 & 0.793 & 0.816 & 0.881 & & 0.862 & 0.842 & 0.855 & 0.790 \\
\bottomrule
\end{tabular}}
\label{tab:cone_strength}
\end{table}

Figure~\ref{fig:geometry} visualizes the joint embedding geometry for a representative subset of datasets and backbones. We observe that across all models, both image and text representations concentrate into narrow angular regions of the hypersphere, confirming strong modality-specific cone structure. Due to space constraints, we present the full quantitative analysis of this phenomenon in Table~\ref{tab:cone_strength}, reporting the Mean Resultant Length ($R$) for all studied datasets. The results confirm that anisotropy is amplified in medical datasets, where reduced semantic diversity yields tighter cones and consistently higher $R$ values compared to natural image datasets.

Despite contrastive training, we observe a persistent centroid displacement between modalities, with the modality gap largely aligned with the dominant principal direction. Medical-domain backbones (BioMedCLIP, MedSigLIP) exhibit smaller cross-modal separation than generalist models, but at the cost of increased cone concentration, suggesting that medical pretraining reshapes but does not eliminate modality misalignment.

Figure~\ref{fig:lambda} evaluates the downstream impact of controlled gap reduction via the alignment strength $\lambda$. Moderate closure of the modality gap consistently improves linear-probe AUC across datasets, indicating that excessive cross-modal displacement is detrimental for transfer. However, performance typically saturates---and in some cases slightly degrades---as $\lambda\!\rightarrow\!1$, particularly for medical-domain models on BRSET and mBRSET. This highlights that full collapse of the gap can remove modality-specific information, supporting the view that the optimal geometry corresponds to an intermediate, task-dependent separation rather than perfect overlap. Generalist models benefit more strongly from post-hoc alignment on medical data, consistent with larger initial misalignment under domain shift, while medical-tuned models show smaller but stable gains.

These results suggest that the modality gap is not merely a defect to minimize, but a controllable geometric degree of freedom: reducing excessive separation improves multimodal transfer, yet preserving some modality-specific structure remains critical for maintaining clinically relevant information.

\begin{figure}[H]
    \centering
    \includegraphics[width=\linewidth]{Figures/Fig2.pdf}
    \caption{Downstream AUC as a function of alignment strength $\lambda$ across datasets and backbones. Moderate reduction of the modality gap consistently improves performance, while excessive closure can lead to saturation or degradation, particularly for medical-domain models. Error bars denote standard deviation across seeds.}
    \label{fig:lambda}
\end{figure}















% =========================
% NOTES:
% - Keep your original \cite{3.28} (Mind the Gap) and \cite{3.25} (CLIP).
% - Add BibTeX entries for SigLIP and any medical variants you cite, e.g.:
%   \cite{siglip,biomedclip,medsiglip}
% =========================


%
% the environments 'definition', 'lemma', 'proposition', 'corollary',
% 'remark', and 'example' are defined in the LLNCS documentclass as well.
%


 %% removed for anonymized MICCAI submission.
    
    % The following acknowledgement and disclaimer sections can be removed for the double-blind review process.  If and when your paper is accepted, reinsert the acknowledgement and the disclaimer clause in your final camera-ready version.
    % IF you opted to include the acknowledgement and disclaimer sections, they will count towards the 8-page limit.

\begin{credits}
\subsubsection{\ackname} A bold run-in heading in small font size at the end of the paper is
used for general acknowledgments, for example: This study was funded
by X (grant number Y).

\subsubsection{\discintname}
It is now necessary to declare any competing interests or to specifically
state that the authors have no competing interests. Please place the
statement with a bold run-in heading in small font size beneath the
(optional) acknowledgments\footnote{If EquinOCS, our proceedings submission
system, is used, then the disclaimer can be provided directly in the system.},
for example: The authors have no competing interests to declare that are
relevant to the content of this article. Or: Author A has received research
grants from Company W. Author B has received a speaker honorarium from
Company X and owns stock in Company Y. Author C is a member of committee Z.
\end{credits}
\bibliographystyle{splncs04}
\bibliography{references} 
\end{document}